% !TEX root = ../main.tex

\begin{innovations}

本论文研究了紧 Kähler 流形上带有除子奇性的向量丛及其推广对象的几何结构，主要关注抛物向量丛与 Hermitian--Einstein 度量之间的关系，并在此基础上讨论了 Kobayashi--Hitchin 对应在更一般情形下的成立问题。围绕这一主题，本文在现有理论基础上进行了若干推广和改进，主要创新性体现在以下几个方面。

首先，本文在传统抛物向量丛理论的框架下，将研究对象推广到饱和自反抛物层，并在适当的稳定性条件下建立了其与 Hermitian--Einstein 度量之间的对应关系。已有的 Kobayashi--Hitchin 对应主要针对光滑情形或向量丛情形，而对于存在奇性结构的自反层情形，相关理论仍不完善。本文通过对奇性解析结构的细致分析，构造了适合自反抛物层的解析框架，从而将经典对应推广到更一般的几何对象。

其次，本文结合 Hironaka 奇点解析理论，对带有奇性结构的抛物层进行适当的解析化处理，并在解析化空间上建立了 Hermitian--Einstein 度量存在性的分析方法。通过在解析化后的空间上构造适当的度量和连接，并利用椭圆型偏微分方程的估计方法，本文给出了 Hermitian--Einstein 方程在奇性背景下的可解性分析，从而为建立对应关系提供了关键的解析工具。

再次，本文在上述结构基础上进一步研究了相关的 Chern 类与稳定性之间的关系，得到了一类 Bogomolov--Gieseker 型不等式。这一不等式为抛物自反层的稳定性提供了必要的数值条件，同时也为研究其模空间结构提供了重要的几何约束。

综上所述，本文在抛物几何与 Hermitian--Einstein 度量理论的交叉研究中，对带奇性几何对象的 Kobayashi--Hitchin 对应进行了系统研究，并在解析方法和几何结构方面给出了一定的推广与补充，为进一步研究带奇性结构的复几何对象提供了新的思路和工具。


\end{innovations}


% \begin{innovations}
%
% 工程博士学位论文应主要聚焦工程实践和应用研究，须体现工程性、创新性、实践性、应用性特征，
% 体现学位申请人在专业领域掌握坚实全面的基础理论和系统深入的专门知识，
% 具有独立承担专业实践工作的能力，在专业实践领域做出创新性成果，
% 对推动本专业领域知识和技术的发展作出重要贡献。
%
% 故在此须说明本学位论文的创新性及应用性，确保符合工程博士学位论文要求，字数在800字以内。
%
% \end{innovations}
